\section{Lecture 4: Ext-Computations in Condensed Abelian Groups}
\label{lec:4}

Scholze's third lecture is devoted entirely to \emph{computing} $\Ext$-groups in
light condensed abelian groups. The point of departure is that $\CondAbl$ is a
very convenient framework for homological algebra: it is a nice abelian
category, one can form its derived category without any fuss, and --- since
topological spaces embed fully faithfully into condensed sets, and (under mild
hypotheses) topological abelian groups embed fully faithfully into condensed
abelian groups --- all the usual objects of functional analysis now live in a
setting with a well-behaved $\Ext$. For the theory to be useful these
$\Ext$-groups must have interpretable answers, and the bulk of the lecture is
about establishing that they do. There are two main computations: cohomology of
a topological space against a discrete coefficient group recovers sheaf (hence,
for CW complexes, singular) cohomology; and the $\Ext$-groups between locally
compact abelian groups recover the classical Yoneda-$\Ext$ of topological
abelian groups. The lecture closes with the key ``solidity'' computation
$\Ext^\ast(\prod_{\mathbb{N}}\mathbb{Z},\mathbb{Z})$ and a striking
set-theoretic theorem about when products turn into direct sums.

Much of this and the following lectures follows Scholze's first course on
condensed mathematics; the relevant reference is the \emph{Condensed
Mathematics} lecture notes \cite{clausen2019condensed}.

\subsection{Condensed cohomology recovers singular cohomology}

Recall from \cref{lec:3} the definition of condensed cohomology: for a condensed
set $X$ and an abelian group $M$ (for now discrete, more generally a condensed
abelian group), one sets
\[
H^\ast(X,M):=\Ext^\ast_{\CondAbl}\bigl(\mathbb{Z}[\underline{X}],\,M\bigr),
\]
the $\Ext$ computed in light condensed abelian groups between the free condensed
abelian group on $X$ and $M$. In any topos there is such an internal notion of
cohomology --- $\Ext$ in the category of abelian sheaves --- and the content of
the theorem stated at the end of \cref{lec:3} (\cref{thm:3:singular}) is that in
the condensed world it recovers singular cohomology.

\begin{theorem}[Singular cohomology, recovered]
\label{thm:4:singular-recovery}
Let $X$ be a CW complex and $M$ an abelian group. Then
\begin{equation}
\label{eq:4:singular-recovery}
\Ext^{i}_{\CondAbl}\bigl(\mathbb{Z}[\underline{X}],\,M\bigr)
\;\cong\;
H^{i}_{\mathrm{sing}}(X,M).
\end{equation}
\end{theorem}

The rest of the subsection sketches the proof, which proceeds through sheaf
cohomology. Two reductions get us from a general CW complex to a compact space.

\begin{proof}[Proof sketch, first reductions]
A CW complex is an increasing union $X=\bigcup_i X_i$ of compact (and metrizable)
Hausdorff subspaces $X_i$. Both sides of \eqref{eq:4:singular-recovery} take
filtered colimits in $X$ to derived limits, so --- with a slightly sharper
statement comparing the full complexes rather than the individual cohomology
groups --- one reduces to the case where $X$ is a compact (CW) space. The point
is then that there is a more general statement, valid for compact Hausdorff
spaces that need not be CW complexes at all.
\end{proof}

\begin{theorem}[Compact Hausdorff case: sheaf cohomology]
\label{thm:4:compact-sheaf}
Let $X$ be a metrizable compact Hausdorff space and $M$ an abelian group. Then
\begin{equation}
\label{eq:4:compact-sheaf}
\Ext^{i}_{\CondAbl}\bigl(\mathbb{Z}[\underline{X}],\,M\bigr)
\;=\;
H^{i}_{\mathrm{sheaf}}(X,M),
\end{equation}
the sheaf cohomology of $X$ with coefficients in the constant sheaf $M$.
\end{theorem}

Here $H^i_{\mathrm{sheaf}}$ is ordinary sheaf cohomology on the topological space
$X$: one has the global-sections functor
\begin{equation}
\label{eq:4:global-sections}
\Gamma(X,-)\colon \mathrm{Shv}(X,\mathrm{Ab})\longrightarrow\mathrm{Ab}
\end{equation}
on the category of sheaves of abelian groups on $X$, and its right derived
functors $H^i_{\mathrm{sheaf}}(X,-)$, applied to the constant sheaf $M$.

It is classical that for a CW complex sheaf and singular cohomology agree ---
$H^i_{\mathrm{sheaf}}$ restricted to CW complexes satisfies the
Eilenberg--Steenrod axioms, hence coincides with $H^i_{\mathrm{sing}}$ --- so
\cref{thm:4:compact-sheaf} implies \cref{thm:4:singular-recovery}. In general,
however, the two differ, and it is sheaf cohomology (equivalently, the
$\Ext$-cohomology) that is the correct invariant.

\begin{example}[Totally disconnected spaces: sheaf vs.\ singular]
\label{ex:4:totally-disconnected}
Let $X$ be a totally disconnected compact Hausdorff space --- i.e.\ a profinite
set. Then $\Gamma(X,-)$ is already exact, so
\begin{equation}
\label{eq:4:sheaf-totally-disc}
H^{i}_{\mathrm{sheaf}}(X,M)=
\begin{cases}
\{\text{locally constant maps }X\to M\} & i=0,\\
0 & i>0.
\end{cases}
\end{equation}
By contrast the singular cohomology is
\begin{equation}
\label{eq:4:sing-totally-disc}
H^{i}_{\mathrm{sing}}(X,M)=
\begin{cases}
\{\text{\underline{all} maps }X\to M\} & i=0,\\
0 & i>0.
\end{cases}
\end{equation}
Indeed the singular cochain complex is built from maps of simplices into $X$, and
from the point of view of simplices there is no way to tell $X$ apart from a
discrete set of points: any simplex is connected, so it factors through a single
point of $X$. Thus singular cohomology sees nothing of the topology of $X$,
whereas sheaf cohomology --- returning the locally constant maps in degree $0$
--- does. This is exactly why in general one should use sheaf cohomology (the
same as the $\Ext$-cohomology), not singular cohomology.
\end{example}

\begin{remark}[Locally contractible and paracompact hypotheses]
\label{rmk:4:locally-contractible}
Scholze indicated the range of validity somewhat tentatively. He expects that
for the $\Ext$-cohomology, crossing $X$ with an interval does not change the
answer, though he was not completely sure. For a \emph{locally contractible}
space one can compare with sheaf cohomology, and if in addition the space is
\emph{paracompact} one can compare sheaf cohomology with singular cohomology (by
the usual results). The relevant notion of locally contractible is the ``funny''
one: not that every point has a basis of contractible open neighborhoods, but
that for any point and any neighborhood there is a smaller neighborhood that can
be contracted \emph{inside the larger one}. With this hypothesis everything one
knows for CW complexes extends. Scholze stressed that these better compactness
assumptions only matter at the level of the intermediate sheaf cohomology: if one
goes all the way back to the $\Ext$-groups themselves, the paracompactness issue
probably disappears again. He also emphasized that he does \emph{not} want to
work with topological spaces up to weak homotopy equivalence --- otherwise a
totally disconnected space would become a discrete set of points and could not be
treated at all; the interest is genuinely in the actual topological space.
\end{remark}

\subsection{The sheaf-cohomology upgrade via a geometric morphism}

\Cref{thm:4:compact-sheaf} can be upgraded from constant coefficients to
arbitrary abelian sheaves, and the cleanest way to see it is through a geometric
morphism of topoi. Fix a metrizable compact Hausdorff space $X$. There are two
sites over $X$:
\begin{equation}
\label{eq:4:two-sites}
\CondSetl_{/\underline{X}}
=\bigl(\ProFin_{\mathbb{N}}(\mathrm{Fin})_{/\underline{X}}\bigr)^{\sim}
\;\xrightarrow{\ \lambda\ }\;
\mathrm{Op}(X)^{\sim}=\mathrm{Shv}(X),
\end{equation}
on the left the topos of sheaves on light profinite sets equipped with a map to
$\underline{X}$ (light condensed sets over $X$), on the right the topos of
sheaves on the poset $\mathrm{Op}(X)$ of open subsets of $X$ --- ordinary
sheaves on the topological space. This $\lambda$ is a geometric morphism: any
sheaf on $X$ pulls back to a light condensed set over $X$, and to define the
pullback it suffices to define it on generators, by
\begin{equation}
\label{eq:4:lambda-star-open}
\lambda^{\ast}\underline{U}=\underline{U}\quad(\to\underline{X}),\qquad
U\subseteq X\ \text{open}.
\end{equation}
Passing to abelian sheaves and derived categories gives
\begin{equation}
\label{eq:4:lambda-star-derived}
\lambda^{\ast}\colon
D\bigl(\mathrm{Shv}(X,\mathrm{Ab})\bigr)\longrightarrow
D\bigl(\mathrm{Ab}(\CondSetl_{/\underline{X}})\bigr).
\end{equation}

\begin{theorem}[Full faithfulness of $\lambda^\ast$]
\label{thm:4:lambda-fully-faithful}
Let $X$ be metrizable compact Hausdorff. On bounded-below objects, $\lambda^\ast$
is fully faithful:
\[
\lambda^{\ast}\colon D^{+}\bigl(\mathrm{Shv}(X,\mathrm{Ab})\bigr)
\hookrightarrow D^{+}\bigl(\mathrm{Ab}(\CondSetl_{/\underline{X}})\bigr).
\]
Consequently, for every $\mathcal{F}\in\mathrm{Shv}(X,\mathrm{Ab})$,
\begin{equation}
\label{eq:4:cohomology-comparison}
H^{i}\bigl(\CondSetl_{/\underline{X}},\,\lambda^{\ast}\mathcal{F}\bigr)
\;\cong\;
H^{i}_{\mathrm{sheaf}}(X,\mathcal{F}).
\end{equation}
Applying this to $\mathcal{F}$ the constant sheaf $M$ recovers
\cref{thm:4:compact-sheaf}.
\end{theorem}

\begin{remark}[Robustness]
\label{rmk:4:robustness}
This is the ``fancy'' statement Scholze alluded to: it is not merely that a
single very specific sheaf (the constant one) has the same cohomology on both
sides, but that the whole bounded-below derived category of sheaves on $X$
embeds fully faithfully into the condensed side, and cohomology of \emph{any}
abelian sheaf is computed either way. In particular the argument works for
non-constant sheaves.
\end{remark}

\begin{proof}[Proof sketch]
By adjunction, full faithfulness of $\lambda^\ast$ is equivalent to the unit map
being an isomorphism: for all $A\in D^{+}(\mathrm{Shv}(X,\mathrm{Ab}))$,
\begin{equation}
\label{eq:4:unit-iso}
A\;\xrightarrow{\ \sim\ }\;R\lambda_{\ast}\,\lambda^{\ast}A.
\end{equation}
This is a map in $D^{+}(\mathrm{Shv}(X,\mathrm{Ab}))$, and whether it is an
isomorphism can be checked on stalks at points $x\in X$:
\begin{equation}
\label{eq:4:stalk-check}
A_{x}\;\xrightarrow{\ \sim\ }\;\bigl(R\lambda_{\ast}\,\lambda^{\ast}A\bigr)_{x}.
\end{equation}
The key point is a \emph{base-change property}: taking the stalk at $x$ commutes
with $R\lambda_{\ast}$. This is where the topos-theoretic nature of
$\underline{X}$ is used. One invokes the general results that cohomology
commutes with filtered colimits for \emph{coherent} topoi, coherence meaning
quasicompact-plus-quasiseparated; the key geometric input is precisely that
$\underline{X}$ is qcqs (\cref{prop:3:qcqs}) --- the same fact, recalled from
\cref{lec:3}, that to compute the cohomology of e.g.\ the interval one needs a
surjection onto it from one of the generating objects of the site, a Cantor set.
Pulling the stalk functor inside $R\lambda_\ast$ then reduces the whole
statement to the case where $X=\{x\}$ is a point. Over a point the assertion is
clear: condensed abelian groups sit fully faithfully, and $\mathbb{Z}$ (as a
condensed abelian group) has the two required cohomology objects, so the unit is
an isomorphism.
\end{proof}

\begin{remark}[Boundedness below and closed neighborhoods]
\label{rmk:4:bounded-below}
The restriction to $D^{+}$ is genuine. Without a bound to the left there are
issues with how cohomology interacts with the passage to limits; there are
several versions of the unbounded derived category (the naive derived category
of sheaves, and its left completion), and to retain full faithfulness one should
use the left completion. For a CW complex, or in finite cohomological dimension,
this makes no difference. When executing the stalk argument one should moreover
regard a point $x$ as the limit of its \emph{closed} neighborhoods rather than
its open ones: in a compact Hausdorff space the closed and open neighborhoods are
both cofinal, but pulling back along a closed neighborhood keeps one within
profinite (qcqs) sets, which is what the coherence argument needs.
\end{remark}

\subsection{A concrete proof by hypercover resolutions}

Scholze also indicated the more hands-on proof of \cref{thm:4:compact-sheaf},
which is instructive and breaks into two steps. Concretely, for $X$ metrizable
compact Hausdorff one wants to compute $\Ext^{i}(\mathbb{Z}[\underline{X}],
\mathbb{Z})$ by finding a projective --- or at least $\Ext^i(-,\mathbb{Z})$-acyclic
--- resolution of $\mathbb{Z}[\underline{X}]$.

\subsubsection*{Step 1: totally disconnected \texorpdfstring{$X$}{X}}

\begin{proposition}[Acyclicity for profinite sets]
\label{prop:4:profinite-acyclic}
Let $S\in\ProFin_{\mathbb{N}}(\mathrm{Fin})$ be a light profinite set. Then
\begin{equation}
\label{eq:4:profinite-ext}
\Ext^{i}\bigl(\mathbb{Z}[S],\,\mathbb{Z}\bigr)=
\begin{cases}
\Cont(S,\mathbb{Z}) & i=0,\\
0 & i>0.
\end{cases}
\end{equation}
\end{proposition}

For $S=\mathbb{N}\cup\{\infty\}$ this is immediate from
\cref{thm:3:three-properties}(3): $\mathbb{Z}[S]$ is (internally) projective. For
general $S$ one must genuinely resolve, and the point is the following. Any
higher $\Ext$-class against $\mathbb{Z}[S]$ can be split after passing to a
suitable cover of $S$; iterating produces a hypercover. Recall that for a
hypercover
% board 18 @ 00:39:52 — /home/deven/documents/coding/vms/gpu/home/notetaker/output/peter-scholze-424-analytic-stacks/boards/board-18.jpg
\[
\begin{tikzcd}
\cdots \arrow[r, shift left=2] \arrow[r] \arrow[r, shift right=2]
& S_2 \arrow[r, shift left=1.5] \arrow[r] \arrow[r, shift right=1.5]
& S_1 \arrow[r, shift left=1] \arrow[r, shift right=1]
& S_0 \arrow[r]
& S
\end{tikzcd}
\]
of light profinite sets --- meaning $S_0\twoheadrightarrow S$, then
$S_1\twoheadrightarrow S_0\times_S S_0$, and $S_{n+1}$ surjecting onto the
appropriate matching object at each stage --- the associated augmented chain
complex of free groups
\begin{equation}
\label{eq:4:hypercover-chain}
\cdots\to\mathbb{Z}[S_2]\to\mathbb{Z}[S_1]\to\mathbb{Z}[S_0]\to\mathbb{Z}[S]\to0
\end{equation}
is exact --- this is a completely general fact about hypercovers in any site.

\begin{proof}[Proof of \cref{prop:4:profinite-acyclic}]
Given the exact resolution \eqref{eq:4:hypercover-chain}, one must show that
applying $\Cont(-,\mathbb{Z})=\underline{\Hom}(\mathbb{Z}[-],\mathbb{Z})$
keeps it exact, i.e.\ that
\begin{equation}
\label{eq:4:cont-complex}
0\to\Cont(S,\mathbb{Z})\to\Cont(S_0,\mathbb{Z})\to\Cont(S_1,\mathbb{Z})\to\cdots
\end{equation}
is exact. There are two ways. \emph{Either} treat every term as the global
sections of a sheaf on $S$: push everything forward to $S$, so that it suffices
to prove exactness of sheaves on $S$, which is checked on stalks, and over a
point a cover is split so the statement is elementary (this is the homological
descent of SGA~4, via the proper base change theorem for ``proper'' = universally
closed separated maps of compact Hausdorff spaces). \emph{Or} --- the argument
Scholze used in the lecture notes --- write the hypercover of profinite sets as a
cofiltered limit of hypercovers of \emph{finite} sets by finite sets; for finite
sets every hypercover is split and the exactness is obvious, and the exactness
passes to the filtered colimit. It takes a little unravelling to see the
hypercover can always be so written, but it reduces to the finite case where the
claim is clear.
\end{proof}

\subsubsection*{Step 2: general \texorpdfstring{$X$}{X}}

Now let $X$ be a general metrizable compact Hausdorff space. We resolve
$\mathbb{Z}[\underline{X}]$ by $\mathbb{Z}[S]$'s with $S$ profinite. Using that
every such $X$ receives a surjection from a Cantor set (indeed from a light
profinite set), pick $S_0\twoheadrightarrow X$ and form the \v{C}ech nerve, the
simplicial light profinite set with
\begin{equation}
\label{eq:4:cech-nerve}
S_i=\underbrace{S_0\times_X\cdots\times_X S_0}_{i+1},
\qquad
S_1=S_0\times_X S_0,\quad S_2=S_0\times_X S_0\times_X S_0,\ \dots
\end{equation}
Because $S_0$ is already totally disconnected, each fibre product $S_i$ is a
closed subspace of a product of totally disconnected spaces, hence again totally
disconnected --- so no further resolution is needed. The \v{C}ech nerve
$S_\bullet\to\underline{X}$ is a hypercover, and by the same general principle
its augmented chain complex is a resolution
\begin{equation}
\label{eq:4:cech-resolution}
\cdots\to\mathbb{Z}[S_2]\to\mathbb{Z}[S_1]\to\mathbb{Z}[S_0]\to
\mathbb{Z}[\underline{X}]\to0
\end{equation}
whose terms $\mathbb{Z}[S_i]$ are $\Ext^i(-,\mathbb{Z})$-acyclic by
\cref{prop:4:profinite-acyclic}. Hence $\Ext^{i}(\mathbb{Z}[\underline{X}],
\mathbb{Z})$ is computed by the complex
\begin{equation}
\label{eq:4:cech-cont-complex}
0\to\Cont(S_0,\mathbb{Z})\to\Cont(S_1,\mathbb{Z})\to\Cont(S_2,\mathbb{Z})\to\cdots,
\end{equation}
an admittedly awkward-looking formula for something like singular cohomology: one
covers $X$ by a Cantor set, forms all the fibre products, takes locally constant
$\mathbb{Z}$-valued functions everywhere, and the claim is that the resulting
complex computes the right thing. That it does is proved exactly as before ---
treat the terms as global sections of sheaves on $X$ and check on stalks that
\eqref{eq:4:cech-cont-complex} resolves the constant sheaf $\underline{\mathbb{Z}}$
on $X$.

\subsection{Locally compact abelian groups}

The second main computation concerns the classical world of locally compact
abelian groups, which embeds fully faithfully into condensed abelian groups and
whose $\Ext$-groups one can therefore ask about.

\begin{notation}[Metrizable LCA groups]
\label{not:4:lca}
Write $\LCA_m$ for the category of metrizable locally compact abelian groups.
Examples include discrete abelian groups, $\mathbb{R}$, the circle
$\mathbb{R}/\mathbb{Z}$, the $p$-adic integers $\mathbb{Z}_p$, and the adeles
\begin{equation}
\label{eq:4:adeles}
\mathbb{A}=\widehat{\mathbb{Z}}\otimes\mathbb{Q}\times\mathbb{R}.
\end{equation}
\end{notation}

The adeles illustrate the general structure: $\mathbb{Q}$ sits inside $\mathbb{A}$
as a discrete subgroup with compact quotient,
\begin{equation}
\label{eq:4:adele-ses}
0\to\mathbb{Q}\to\mathbb{A}\to\mathbb{A}/\mathbb{Q}\to0,
\end{equation}
with $\mathbb{Q}$ discrete and $\mathbb{A}/\mathbb{Q}$ compact. In general every
$A\in\LCA_m$ admits a filtration with three pieces:
\begin{equation}
\label{eq:4:lca-filtration}
\text{a discrete group,}\quad
\text{a finite-dimensional real vector space,}\quad
\text{a compact metrizable abelian group.}
\end{equation}

\begin{theorem}[Condensed $\Ext$ of LCA groups]
\label{thm:4:lca-ext}
For $A,B\in\LCA_m$,
\begin{equation}
\label{eq:4:lca-ext}
\Ext^{i}_{\CondAbl}\bigl(\underline{A},\underline{B}\bigr)\cong
\begin{cases}
\Hom_{\LCA}(A,B) & i=0,\\
\Ext^{1}_{\LCA}(A,B) & i=1,\\
0 & i\geq2,
\end{cases}
\end{equation}
where $\Ext^{1}_{\LCA}$ is the classical Yoneda $\Ext$ of extensions of $A$ by
$B$ within locally compact abelian groups (up to the appropriate notion of
equivalence).
\end{theorem}

Full faithfulness already forces the $i=0$ term to be the ordinary
(topological) homomorphisms. What is not a priori clear is that no exotic higher
$\Ext$-classes appear: the condensed $\Ext^1$ turns out to be exactly the
classical Yoneda $\Ext^1$, and everything above degree $1$ vanishes --- matching
the classical fact that $\LCA$ is an exact category with $\Ext^i_{\LCA}=0$ for
$i\geq2$. (Scholze noted this is special: there \emph{are} pairs of condensed
abelian groups with surprising higher $\Ext$, and the solidity computation below
is one such.)

\begin{example}[Pontryagin duality]
\label{ex:4:pontryagin}
For any $A\in\LCA_m$,
\begin{equation}
\label{eq:4:pontryagin}
\Ext^{i}\bigl(\underline{A},\underline{\mathbb{R}/\mathbb{Z}}\bigr)=
\begin{cases}
A^{\vee} & i=0,\\
0 & i>0,
\end{cases}
\end{equation}
where $A^{\vee}=\Hom(A,\mathbb{R}/\mathbb{Z})$ is the Pontryagin dual. Thus
$A\mapsto A^{\vee}$ is exact on locally compact abelian groups, sitting in degree
$0$ with nothing in higher degrees --- the condensed reflection of Pontryagin
duality.
\end{example}

\begin{example}[No maps from the reals to the integers]
\label{ex:4:R-Z}
\begin{equation}
\label{eq:4:R-Z}
\Ext^{i}\bigl(\underline{\mathbb{R}},\mathbb{Z}\bigr)=0\qquad\text{for all }i\geq0.
\end{equation}
The intuition is that $\mathbb{R}$ is connected, so it can never map to the
discrete group $\mathbb{Z}$ (giving $\Ext^0=0$), and all higher groups vanish as
well. This is a computation one really has to do; the vanishing is proved below.
\end{example}

\subsection{The Breen--Deligne resolution}

To compute anything one needs something close to a projective resolution. After
a d\'evissage reducing $A$ and $B$ to the three basic pieces of
\eqref{eq:4:lca-filtration} (mapping to a discrete group is easy; the
finite-dimensional real vector spaces are handled separately), one may assume $A$
is a compact metrizable group. The engine is a functorial resolution valid for
\emph{all} abelian groups.

\begin{theorem}[Breen--Deligne resolution]
\label{thm:4:breen-deligne}
There is a resolution, functorial in the abelian group $M$,
\begin{equation}
\label{eq:4:breen-deligne}
\cdots\to\mathbb{Z}[M^{n_i}]\to\cdots\to\mathbb{Z}[M^{2}]\to\mathbb{Z}[M]\to M\to0,
\end{equation}
in which every term is a single free abelian group $\mathbb{Z}[M^{n_i}]$ on a
power of $M$ and all differentials are given by \emph{universal formulas}. The
augmentation is $[m]\mapsto m$ and the first differential is
\begin{equation}
\label{eq:4:bd-first-differential}
[(a,b)]\longmapsto [a]+[b]-[a+b].
\end{equation}
\end{theorem}

\begin{remark}[Provenance and (in)explicitness]
\label{rmk:4:breen-deligne-provenance}
This resolution was used substantially by Breen in an algebraic-geometry setting,
but the statement in the form needed here was never put in the literature by him;
there is an unpublished letter of Deligne to Breen proving the result. Two
features are striking. First, no \emph{explicit} such resolution is known: the
differentials exist and are natural, but one cannot write them down, and Scholze
noted (with mild surprise) that the proof uses a little stable homotopy theory,
in some incarnations the finiteness of the stable homotopy groups of spheres ---
one has to kill something like the stable stems, and there is no explicit way to
do so. Second, one can arrange a \emph{single} free group $\mathbb{Z}[M^{n_i}]$
in each degree rather than a finite direct sum: any such direct sum can itself be
covered by one free group of the same kind, so one chooses one term per degree.
\end{remark}

Because \eqref{eq:4:breen-deligne} is given by universal formulas, functoriality
makes it work for abelian sheaves on any site. Applied to a condensed abelian
group $A$ it gives
\begin{equation}
\label{eq:4:breen-deligne-A}
\cdots\to\mathbb{Z}[\underline{A}^{n_i}]\to\cdots\to\mathbb{Z}[\underline{A}]\to
\underline{A}\to0,
\end{equation}
which reduces the computation of $\Ext^{i}(\underline{A},-)$ to that of
$\Ext^{i}(\mathbb{Z}[\underline{A}^{n}],-)$ --- and these are exactly the
cohomology computations of the previous subsections, since $\underline{A}^{n}$ is
a nice space. One more input is needed for real coefficients.

\begin{proposition}[Cohomology with real coefficients]
\label{prop:4:real-coefficients}
For $X$ metrizable compact Hausdorff,
\begin{equation}
\label{eq:4:real-coefficients}
\Ext^{i}_{\CondAbl}\bigl(\mathbb{Z}[\underline{X}],\underline{\mathbb{R}}\bigr)=
\begin{cases}
\Cont(X,\mathbb{R}) & i=0,\\
0 & i>0.
\end{cases}
\end{equation}
The same holds with $\mathbb{R}$ replaced by any Banach space $V$, but this uses
the \emph{local convexity} of $V$: the proof runs through partition-of-unity
arguments, which require the target to have a basis of convex neighborhoods.
\end{proposition}

\begin{remark}[Non-locally-convex coefficients]
\label{rmk:4:non-convex-preview}
As a preview: when real vector spaces are treated inside the analytic theory
later, one is forced to consider \emph{non-locally-convex} vector spaces, for
which \cref{prop:4:real-coefficients} can fail --- suspension need not behave
well on all compact Hausdorff spaces. There one really has to resolve by totally
disconnected objects instead.
\end{remark}

\begin{example}[$\Ext^{i}(\mathbb{R},\mathbb{Z})=0$, worked out]
\label{ex:4:R-Z-worked}
Resolve $\underline{\mathbb{R}}$ by the Breen--Deligne resolution
\eqref{eq:4:breen-deligne-A}, with terms $\mathbb{Z}[\underline{\mathbb{R}}^{n}]$.
Since $\mathbb{R}^{n}$ is contractible,
\begin{equation}
\label{eq:4:ZRn-Z}
\Ext^{i}\bigl(\mathbb{Z}[\underline{\mathbb{R}}^{n}],\mathbb{Z}\bigr)
=H^{i}_{\mathrm{sing}}(\mathbb{R}^{n},\mathbb{Z})=
\begin{cases}
\mathbb{Z} & i=0,\\
0 & i>0,
\end{cases}
\end{equation}
so each term contributes exactly one copy of $\mathbb{Z}$, and one must see that
these assemble to $0$. The trick is to run the \emph{same} resolution on the zero
group, which resolves $0$; comparing the two augmented complexes column by
column,
% board 28 @ 01:20:04 — /home/deven/documents/coding/vms/gpu/home/notetaker/output/peter-scholze-424-analytic-stacks/boards/board-28.jpg
\[
\begin{tikzcd}
\cdots \arrow[r] & \mathbb{Z}[0] \arrow[r] \arrow[d] & \cdots \arrow[r] & \mathbb{Z}[0] \arrow[r] \arrow[d] & \mathbb{Z}[0] \arrow[r] \arrow[d] & 0 \arrow[d] \\
\cdots \arrow[r] & \mathbb{Z}[\underline{\mathbb{R}}^{n}] \arrow[r] & \cdots \arrow[r] & \mathbb{Z}[\underline{\mathbb{R}}^{2}] \arrow[r] & \mathbb{Z}[\underline{\mathbb{R}}] \arrow[r] & \underline{\mathbb{R}}
\end{tikzcd}
\]
the map $\mathbb{Z}[0^{n}]=\mathbb{Z}[0]\to\mathbb{Z}[\underline{\mathbb{R}}^{n}]$
induces an isomorphism on each $\Ext^{\ast}(-,\mathbb{Z})$ by \eqref{eq:4:ZRn-Z}.
Hence $\Ext^{\ast}(\underline{\mathbb{R}},\mathbb{Z})$ agrees with
$\Ext^{\ast}$ computed from the top row, which resolves $0$; so it vanishes.
\end{example}

\subsection{The cubical \texorpdfstring{$Q$}{Q}-construction}

One might be uneasy about running explicit computations off an inexplicit
resolution. In fact the inexplicitness never causes trouble --- only the
existence and functoriality of \cref{thm:4:breen-deligne} are used --- but there
is also an \emph{explicit} complex that can be used instead: MacLane's
$Q$-construction, also called the cubical construction, worked out by Commelin in
the course of the formalization effort.

\begin{definition}[The cubical construction]
\label{def:4:cubical}
For an abelian group $M$, let $Q(M)$ be the explicit complex
\begin{equation}
\label{eq:4:Q-complex}
Q(M):\qquad
\cdots\to\mathbb{Z}[M^{8}]\to\mathbb{Z}[M^{4}]\to\cdots\to
\mathbb{Z}[M^{2}]\to\mathbb{Z}[M],
\end{equation}
with $i$-th term $\mathbb{Z}[M^{2^{i}}]$. The first differential is
$[(a,b)]\mapsto[a]+[b]-[a+b]$, and the next is the ``cube'' differential: viewing
$(a,b,c,d)$ as the four corners of a square,
\begin{equation}
\label{eq:4:Q-cube-differential}
[(a,b,c,d)]\longmapsto
[(a,b)]+[(c,d)]-[(a+c,b+d)]-[(a,c)]-[(b,d)]+[(a+b,c+d)],
\end{equation}
subtracting the opposite faces in each of the two directions; higher differentials
arrange the $2^{i}$ arguments on the vertices of an $i$-cube and take the
alternating sum of the pairs of opposite faces.
\end{definition}

\begin{remark}
\label{rmk:4:cube-signs}
Scholze cautioned that he might have the signs in
\eqref{eq:4:Q-cube-differential} slightly wrong; if $d^{2}=0$ fails as written
there is an easy variant that works. Editorially: for the two differentials as
transcribed here the composite does vanish --- expanding
$d_1\circ d_2[(a,b,c,d)]$ via $[(x,y)]\mapsto[x]+[y]-[x+y]$, all nine terms
cancel in pairs, so $d^2=0$ holds at least for this first composite.
\end{remark}

The key property is that $Q$ is, up to torsion corrections, \emph{linear} in
$M$.

\begin{theorem}[Linearity of the cubical construction]
\label{thm:4:cubical-linearity}
There is a natural equivalence
\begin{equation}
\label{eq:4:Q-linearity}
Q(M)\;\simeq\;Q(\mathbb{Z})\otimes^{\mathbb{L}}_{\mathbb{Z}}M,
\end{equation}
a derived base change; if $M$ is torsion-free this reduces to
$H_{i}(Q(M))\cong H_{i}(Q(\mathbb{Z}))\otimes_{\mathbb{Z}}M$, with an additional
$\mathrm{Tor}$ term in general. Moreover
\begin{equation}
\label{eq:4:QZ}
Q(\mathbb{Z})\;\simeq\;\bigoplus_{i\geq0}\bigl(\mathbb{Z}\otimes^{\mathbb{L}}_{\mathbb{S}}
\mathbb{Z}\bigr)^{\oplus 2^{i}}[i],
\end{equation}
the tensor product being over the sphere spectrum $\mathbb{S}$; computing it uses
a little stable homotopy theory. Consequently, for $A\in\CondAbl$ and any $B$,
\begin{equation}
\label{eq:4:Q-vanishing-equiv}
\Ext^{i}(\underline{A},B)=0\ \ \forall i>0
\quad\Longleftrightarrow\quad
\Ext^{i}(Q(\underline{A}),B)=0\ \ \forall i>0.
\end{equation}
\end{theorem}

This linearity in $M$ turns out to be exactly what is needed: any of the
vanishing statements one wants to prove can be put in the form ``$\Ext^{i}=0$ for
all $i\geq i_0$'', and one may then use the explicit complex $Q(A)$ in place of a
resolution of $A$ (it is not literally a resolution, but for this purpose it is
just as good). The whole point is comfort: nothing in the arguments changes, but
one knows the objects involved explicitly, and the shape of the answer ``should
have something to do with'' the stable stems whenever one writes down something
explicit.

\subsection{The solidity computation}

The final computation is the one that seeds the theory of solid abelian groups.
It concerns an object that is no longer locally compact but genuinely large: the
countable product $\prod_{\mathbb{N}}\mathbb{Z}$.

\begin{corollary}[$\Ext$ out of $\prod_{\mathbb{N}}\mathbb{Z}$]
\label{cor:4:solid-generator}
For every discrete abelian group $M$,
\begin{equation}
\label{eq:4:solid-ext}
\Ext^{i}_{\CondAbl}\Bigl(\prod_{\mathbb{N}}\mathbb{Z},\,M\Bigr)=
\begin{cases}
\bigoplus_{\mathbb{N}}M & i=0,\\
0 & i>0.
\end{cases}
\end{equation}
\end{corollary}

The degree-$0$ statement is the classical fact that a homomorphism
$\prod_{\mathbb{N}}\mathbb{Z}\to M$ factors through finitely many coordinates; the
critical new content is that there is \emph{nothing in higher degrees}. This is
very different from the naive classical guess and is exactly what makes
$\prod_{\mathbb{N}}\mathbb{Z}$ a compact projective generator of the full
subcategory
\begin{equation}
\label{eq:4:solid-subcat}
\Solid\subset\CondAbl
\end{equation}
of solid abelian groups (cf.\ \cref{prop:1:Zsolid-structure}, where the products
$\prod_{I}\mathbb{Z}$ appeared as the compact projective generators of the heart
of $D(\Zsolid)$): if one wants $\prod_{\mathbb{N}}\mathbb{Z}$ to be projective,
one had certainly better have all the higher $\Ext$-groups vanish.

\begin{remark}[A warning: $\prod_{\mathbb{N}}\mathbb{Z}$ is a huge colimit]
\label{rmk:4:huge-colimit}
As a condensed set, $\prod_{\mathbb{N}}\mathbb{Z}$ is the filtered union
\begin{equation}
\label{eq:4:prodZ-colimit}
\prod_{\mathbb{N}}\mathbb{Z}=\bigcup_{f\colon\mathbb{N}\to\mathbb{N}}\
\prod_{n\in\mathbb{N}}[-f(n),f(n)],
\end{equation}
over all functions $f\colon\mathbb{N}\to\mathbb{N}$, where each
$\prod_{n}[-f(n),f(n)]$ (a product of finite integer intervals) is a profinite
set. This is a \emph{huge} colimit --- of size the continuum. So the naive
approach, resolving by free groups $\mathbb{Z}[S]$ on profinite sets, is
essentially impossible to execute: mapping out of such a huge colimit turns into
an extremely hard-to-control huge limit.
\end{remark}

\begin{proof}[Proof of \cref{cor:4:solid-generator}]
The trick is to resolve in the ``other'' direction, embedding
$\prod_{\mathbb{N}}\mathbb{Z}$ into a product of copies of $\mathbb{R}$. There is a
short exact sequence
\begin{equation}
\label{eq:4:prod-ses}
0\to\prod_{\mathbb{N}}\mathbb{Z}\to\prod_{\mathbb{N}}\underline{\mathbb{R}}\to
\prod_{\mathbb{N}}\underline{\mathbb{R}/\mathbb{Z}}\to0,
\end{equation}
in which $\prod_{\mathbb{N}}\mathbb{R}/\mathbb{Z}$ is a metrizable compact abelian
group (a countable product of circles). By the compact/LCA computation
(\cref{thm:4:lca-ext} applied to this compact group),
\begin{equation}
\label{eq:4:ext-prod-torus}
\Ext^{i}\Bigl(\prod_{\mathbb{N}}\underline{\mathbb{R}/\mathbb{Z}},\,M\Bigr)=
\begin{cases}
0 & i=0,\\
\bigoplus_{\mathbb{N}}M & i=1,\\
0 & i>1.
\end{cases}
\end{equation}
It remains to show
\begin{equation}
\label{eq:4:ext-prod-R-vanish}
\Ext^{i}\Bigl(\prod_{\mathbb{N}}\underline{\mathbb{R}},\,M\Bigr)=0
\qquad\forall i\geq0.
\end{equation}
Granting \eqref{eq:4:ext-prod-R-vanish}, the long exact sequence obtained by
applying $\Ext^{\ast}(-,M)$ to \eqref{eq:4:prod-ses} gives
$\Ext^{i}(\prod_{\mathbb{N}}\mathbb{Z},M)\cong
\Ext^{i+1}(\prod_{\mathbb{N}}\mathbb{R}/\mathbb{Z},M)$, which by
\eqref{eq:4:ext-prod-torus} is $\bigoplus_{\mathbb{N}}M$ for $i=0$ and $0$ for
$i>0$, as claimed.

For \eqref{eq:4:ext-prod-R-vanish} one uses that $\prod_{\mathbb{N}}\mathbb{R}$
carries a module structure over the condensed ring $\underline{\mathbb{R}}$.
Hence, by adjunction,
\begin{equation}
\label{eq:4:rhom-module}
\RHom\Bigl(\prod_{\mathbb{N}}\underline{\mathbb{R}},\,M\Bigr)
=\RHom_{\underline{\mathbb{R}}}\Bigl(\prod_{\mathbb{N}}\underline{\mathbb{R}},\,
\underbrace{\RHom(\underline{\mathbb{R}},M)}_{=\,0}\Bigr)=0,
\end{equation}
the internal $\RHom(\underline{\mathbb{R}},M)$ vanishing because $\mathbb{R}$ is
connected and $M$ discrete (the generalization of \cref{ex:4:R-Z-worked} to any
discrete $M$). One does not even need to know what
$\RHom(\underline{\mathbb{R}},M)$ is --- only that
$\prod_{\mathbb{N}}\mathbb{R}$ is an $\mathbb{R}$-module and that $\mathbb{R}$
admits no maps to a discrete module.
\end{proof}

\begin{remark}[Hilbert-cube alternative]
\label{rmk:4:hilbert-cube}
There is a second route to \eqref{eq:4:ext-prod-R-vanish}. Running the analogous
operation on $\prod_{\mathbb{N}}\mathbb{R}$ turns each term of the resolution into
a product of intervals --- a Hilbert cube. The comparison
\eqref{eq:4:real-coefficients} between condensed and singular cohomology holds
for the Hilbert cube as well (no higher cohomology), so the argument given for
$\mathbb{R}$ applies verbatim to the Hilbert cube variant.
\end{remark}

\subsection{A set-theoretic theorem}

Scholze closed with a ``fun theorem with a set-theoretic proof''. It asks in what
generality the phenomenon of \cref{cor:4:solid-generator} --- pulling a product
out as a direct sum --- persists, not just for products but for sequential
limits.

\begin{definition}[The assertion $(\ast)$]
\label{def:4:star}
Let $(\ast)$ be the following assertion. For all sequential limits
\begin{equation}
\label{eq:4:sequential-limit}
\cdots\to M_1\to M_0
\end{equation}
of countable discrete abelian groups with surjective transition maps, and all
discrete abelian groups $N$,
\begin{equation}
\label{eq:4:star}
\Ext^{i}\Bigl(\varprojlim_{n}M_n,\,N\Bigr)\;\cong\;
\operatorname*{colim}_{n}\Ext^{i}(M_n,N)\qquad(=0\ \text{for }i\geq2).
\end{equation}
\end{definition}

The right-hand vanishing for $i\geq2$ is automatic, since $\Ext^i(M_n,N)=0$ for
$i\geq2$ between discrete groups; the content is in degrees $0$ and $1$.
Computing such a limit can be resolved by a two-term complex involving products
of the $M_n$, which reduces the limit to a product; resolving the $M_n$ by
countable free abelian groups then reduces $(\ast)$ to the vanishing
\begin{equation}
\label{eq:4:star-reduction}
\Ext^{i}\Bigl(\prod_{\mathbb{N}}\bigoplus_{\mathbb{N}}\mathbb{Z},\,
\bigoplus_{I}\mathbb{Z}\Bigr)=0\qquad\text{for }i>0.
\end{equation}

\begin{remark}[$(\ast)$ fails under CH]
\label{rmk:4:star-CH}
Clausen and Scholze quickly realized that $(\ast)$ \emph{fails} under the
continuum hypothesis $2^{\aleph_0}=\aleph_1$: there one finds
$\Ext^{1}\neq0$. They thought it might be too much to hope for at all --- but it
turned out that set theorists had independently studied precisely the same
question in a different language. Unravelling, the relevant $\Ext$-group is a
huge filtered colimit, over functions $f\colon\mathbb{N}\to\mathbb{N}$, of
products of finite direct sums $\prod_{n}\bigoplus_{m\leq f(n)}\mathbb{Z}$, and
computing it amounts to a huge derived limit --- exactly the higher derived
limits studied in that literature
\cite{bergfalk2021simultaneouslyvanishinghigherderived,bergfalk2021simutaneouslyvanishinghigherderived}.
\end{remark}

\begin{theorem}[Bergfalk--Lambie-Hanson--Hru\v{s}\'ak--Bannister]
\label{thm:4:star-set-theory}
\leavevmode
\begin{enumerate}
\item $(\ast)\implies 2^{\aleph_0}>\aleph_{\omega}$.
\item It is consistent that $(\ast)$ holds and
$2^{\aleph_0}=\aleph_{\omega+1}$; in fact $(\ast)$ holds in any forcing extension
obtained by adjoining $\beth_{\omega}$ many Cohen reals.
\end{enumerate}
\end{theorem}

So $(\ast)$ forces the continuum to be large; it cannot equal $\aleph_{\omega}$
(by K\"onig's theorem), and the first possibility above it,
$2^{\aleph_0}=\aleph_{\omega+1}$, is already consistent with $(\ast)$. The
consistency is realized by Cohen's most basic forcing --- the one he invented to
make the continuum hypothesis false, adjoining new real numbers to a ground model
--- here adjoining $\beth_{\omega}$ many Cohen reals, $\beth_{\omega}$ being in
some sense the minimal number of reals one can adjoin and still have a
chance.

\begin{remark}[Why it is mentioned]
\label{rmk:4:star-usage}
The principle $(\ast)$ is never actually used in the course. It is nonetheless
useful to know it can be arranged: when computing certain things it is easy to
guess what the answer would be if $(\ast)$ held, and the results one genuinely
needs can usually be proved without invoking it. But in some situations one might
want it --- for instance to control $\Ext$-groups out of Fr\'echet-type spaces,
where the expected answer holds under $(\ast)$; in general these $\Ext$-groups
are genuinely set-theory-dependent, their value tied to the order type of the
poset of functions $\mathbb{N}\to\mathbb{N}$ under eventual dominance, which
depends heavily on the model of set theory.
\end{remark}
